% ============================================================
% 【物理骨架 0911｜零参数继承件】＝数学 L1 答案册字体件原样继承（来源
%   工作区/字替对照-0909/靠齐样张-0910/答案册/qp-answ-fonts.tex，md5 13967113…）。
%   与主件 qp-fonts.tex 二选一装载（宏名有重叠），件侧勿同时 \input。
% ============================================================
% qp-answ-fonts.tex —— 字体模块（答案册式样草案）
% 依据：附则《全品基准总表》§字号/字距行＋《全品样张细节清单》§八（v4.3/v4.4 落库）
%       ＋variantF/qp-fonts.tex（H1 落定档，2026-09-10）。
% 纪律：xeCJK 一律 Path= 显式路径（TinyTeX fontconfig 不扫用户级字体目录，已实证）；
%       Path= 值内禁宏（fontspec 对其 detokenize），故路径以字面量书写。
% 复用：黑体族静态实例与楷/仿宋取自 variantF/fonts/（只读引用，不复制、不改动）。
% ============================================================
\setmainfont{Times New Roman}
% 正文＝方正书宋 FZShuSong-Z01S（全品嵌字名同款，MD5 校验一致；只测量不取用见附则§八）
\setCJKmainfont[
  Path=C:/Users/28120/AppData/Local/Microsoft/Windows/Fonts/,
  BoldFont=FZHei-B01S.ttf]{FZShuSong-Z01S.ttf}
% 楷体/仿宋＝方正四款免费字体原件（与全品真身 FZKTK/FZFSK 同源同款）
\setCJKfamilyfont{kai}[Path=C:/提示词/工作区/字替对照-0909/variantF/fonts/]{FZKTK.TTF}
\setCJKfamilyfont{fangsong}[Path=C:/提示词/工作区/字替对照-0909/variantF/fonts/]{FZFSK.TTF}
% ---- 黑体字重阶梯（思源黑体真字重静态实例；全品六族谱替代档，兰亭黑采购遗留） ----
\setCJKfamilyfont{heizhang}[Path=C:/提示词/工作区/字替对照-0909/variantF/fonts/,BoldFont=NSC-w600.ttf]{NSC-w600.ttf}
\setCJKfamilyfont{hejie}[Path=C:/提示词/工作区/字替对照-0909/variantF/fonts/,BoldFont=NSC-w315.ttf]{NSC-w315.ttf}
\setCJKfamilyfont{hekeshi}[Path=C:/提示词/工作区/字替对照-0909/variantF/fonts/,BoldFont=NSC-w600.ttf]{NSC-w600.ttf}
\setCJKfamilyfont{heibf}[Path=C:/提示词/工作区/字替对照-0909/variantF/fonts/,BoldFont=NSC-w600.ttf]{NSC-w600.ttf}
\setCJKfamilyfont{heihao}[Path=C:/提示词/工作区/字替对照-0909/variantF/fonts/,BoldFont=NSC-w700.ttf]{NSC-w500.ttf}
% 数字族走 fontspec 拉丁族（xeCJK 把 ASCII 路由主字体，CJK 族字重对数字失效）；
% Height/Depth 按 Times hhea 口径（asc 0.905／desc 0.212em）裁齐，防行盒顶溢。
\newfontfamily{\numboldjian}[Path=C:/提示词/工作区/字替对照-0909/variantF/fonts/,Height=0.905,Depth=0.212]{NSC-w700.ttf}
\newfontfamily{\numbold}[Path=C:/提示词/工作区/字替对照-0909/variantF/fonts/,Height=0.905,Depth=0.212]{NSC-w850.ttf}
% 裁5（0911 定案·随改）：页码数字族＝NotoSansSC-Medium（w500，同 variantF H3 档；全品竖笔 5–6px/墨高 39px）
\newfontfamily{\numpnum}[Path=C:/提示词/工作区/字替对照-0909/variantF/fonts/]{NotoSansSC-Medium.otf}
% ================= 访问宏 =================
\newcommand{\heiti}{\CJKfamily{heibf}}
\newcommand{\heizhang}{\CJKfamily{heizhang}}
\newcommand{\hejie}{\CJKfamily{hejie}}
\newcommand{\hekeshi}{\CJKfamily{hekeshi}}
\newcommand{\heibf}{\CJKfamily{heibf}}
\newcommand{\heihao}{\CJKfamily{heihao}}
\newcommand{\kaishu}{\CJKfamily{kai}}
\newcommand{\fangsong}{\CJKfamily{fangsong}}
% 字符类：①①、★、数学符号区归 CJK（全品基准总表×variantF 同款）
\xeCJKDeclareCharClass{CJK}{"2460->"24FF, "2600->"26FF, "2200->"22FF, "25A0->"25FF}
\xeCJKDeclareCharClass{Default}{"221A}
% H1 落定档：PunctStyle=banjiao（标点占宽压 0.5em、前后墨空档 15/16px）；
%   CJKglue 0.10em——全品行内字心距 3.99mm＝字号＋0.10em 档（每行 21 汉字×3.986mm＝83.7mm）
\xeCJKsetup{CheckSingle=true,PunctStyle=banjiao,CJKglue={\hskip 0.10em plus 0.05em minus 0.01em}}
